\documentclass{article}
\usepackage{amsmath,amssymb,booktabs,array}
\usepackage[T1]{fontenc}
\usepackage[margin=1in]{geometry}
\usepackage{graphicx}
\usepackage{url}
\setlength{\emergencystretch}{1em}

% Hybrid Upside/dry-MARTINI interface: method, calibration, and evidence limits.

\title{A Hybrid Machine-Learned/Physical Model for Membrane-Protein
Thermodynamics: Coupling the Upside Protein Force Field to a Full-Resolution
dry-MARTINI Membrane}
\author{}
\date{}

\begin{document}
\maketitle

\begin{abstract}
We couple the machine-learned Upside protein force field to a full-resolution
dry-MARTINI membrane without fitting the combined model to membrane-protein data.
All real degrees of freedom use one g-JF force schedule and one system temperature
for both fluctuation--dissipation noise and friction calibration. The preferred
kinetic target is factor-four-corrected DOPC molecular lateral diffusion. A stable
one-step friction scan does not reach that target, so production uses a narrower
fallback: the requested free-particle mobility of a MARTINI bead. Each
membrane-contacting protein N/CA/C carrier is assigned the sum of the frictions of
the lipid beads within the 12~\AA{} interaction range. Regenerated O and
backbone-centre sites are virtual. The backbone-centre reaction is routed with
the MARTINI mass fractions $14/54$, $12/54$, and $12/54$ to persistent N/CA/C;
the $16/54$ share assigned to regenerated O is intentionally disposable under
the Upside integration cycle. This defines the implemented coupling contract,
but not a physical molecular clock. Measured DOPC molecular diffusion
remains orders of magnitude below target, and the contact-count rule is not a
hydrodynamic resistance tensor. Equilibrium observables are potentially usable
because their values do not require a physical trajectory clock, but only after
timestep, ensemble, and replica or enhanced-sampling convergence have been
demonstrated. Present trajectories do not yet pass those gates. In particular,
the HDX adapter correctly reuses the established EX2 analysis, but current HDX
predictions are not quantitative because protection-state populations are
unconverged and a dry membrane supplies no calibrated water-accessibility field.
\end{abstract}

\section{Introduction}

\subsection{The data problem, and two responses}

Machine-learned force fields are powerful exactly where data is abundant. Upside
represents each residue by a few interaction sites, evaluates a learned force
field through pre-tabulated splines, and propagates the chain at a comparatively
large timestep; trained against the large soluble-protein record, it delivers
folding and conformational thermodynamics at low cost. Membrane proteins are the
opposite regime: experimental structures and stabilities are comparatively few,
and the environment (the bilayer) is an active thermodynamic participant that a
protein-only model does not contain.

There are two ways to extend a protein model into the membrane. The first is to
\emph{learn} an implicit membrane term from membrane-protein data---an approach a
previous effort in our group pursued, training an implicit-membrane potential
against on the order of forty membrane proteins. Such a model is only as broad as
its training set: forty proteins is a thin basis for a universal
membrane-protein potential, and a learned implicit term risks encoding the
idiosyncrasies of the training proteins rather than transferable membrane physics.

The second route---ours---adds \emph{no} membrane-protein training data at all.
We keep the machine-learned protein force field where it is strong (protein
internal energetics, learned from abundant soluble data) and supply the membrane
from a \emph{physical}, transferable lipid model (dry-MARTINI), whose parameters
come from lipid thermodynamics, not from protein structures. The protein and the
membrane are then coupled through the \emph{raw} dry-MARTINI cross-interactions,
with a single geometric projection to reconcile the one genuine resolution
mismatch (multi-bead side chain vs.\ single Upside site). Nothing in the
protein model, the membrane model, or their coupling is fit to membrane-protein
data. The power of the combination is precisely that it composes a data-rich
learned model with a transferable physical model, so it can be applied where the
\emph{combined} system has little data---which is the membrane-protein regime, and,
as we note in the outlook, other composite systems such as nanoparticles.

\subsection{What must be true, and what we claim}

For this composition to be useful it must be (i) faithful in the domain it is used
for, (ii) universal (one force field for all proteins; no per-system tuning), and
(iii) affordable at Upside's coarse step. We are deliberately precise about (i):
the hybrid is intended as a \emph{thermodynamic} instrument, not as a source of
absolute membrane-coupled rates. That intent is not itself validation. Every
equilibrium claim still requires timestep convergence, adequate sampling, and a
force field appropriate to the observable. The implemented calibration fixes a
bare-particle mobility and defines a local protein--bilayer friction model; it
does not reproduce DOPC molecular diffusion. Free-energy methods may remove the
need for physical kinetics, but they do not remove the need for correct statics
and converged populations. The same distinction controls the interpretation of
HDX (Sec.~\ref{sec:hdx}).

The contributions are: (1) a full-resolution dry-MARTINI membrane inside Upside
with every pair interaction served from a spline table (Sec.~\ref{sec:membrane});
(2) a shared g-JF force update with one stochastic bath per real degree of freedom
(Sec.~\ref{sec:langevin}); (3) a design-space analysis showing that the tested
single-step hard-core schemes do not reproduce molecular DOPC diffusion and
therefore require an explicitly limited particle-mobility fallback
(Secs.~\ref{sec:designspace},~\ref{sec:timescale}); (4) a validated
area-sampling algorithm and baseline membrane-structure measurements, whose
current-configuration convergence remains to be established
(Sec.~\ref{sec:mechanics}); and (5) a raw-force-field protein--membrane interface
that confines an embedded protein without any added potential
(Sec.~\ref{sec:interface}), enabling conditional equilibrium and HDX analyses
(Sec.~\ref{sec:hdx}).

\section{Theory and Methods}

\subsection{Units, timestep, and per-step decomposition}
\label{sec:units}

Upside uses reduced units: energy $E_\mathrm{up}=2.9150$~kJ/mol, length in
{\AA}ngstr\"om, mass $m_\mathrm{up}=12$~g/mol, and temperature
$1.0\,T_\mathrm{up}=350.59$~K. Protein and lipids share the same numerical step,
$dt=0.009\,t_\mathrm{up}$, used by the standard Upside examples. Each physical
particle is advanced once per step; there is no lipid substep or displacement cap.
There is one system temperature $T$. The workflow passes the same value to the
Upside thermostat and to the friction calibration, so the $k_BT$ used in the
Einstein relation is the same $k_BT$ used in the random impulse. The current
default is $T=0.8\,T_\mathrm{up}=280.47$~K; a different experimental temperature
must be set through this single variable. All native dry-MARTINI parameters (nm, kJ/mol, $e$)
are converted to Upside units once, at spline-table build time; no conversion
constants are baked into the force-field artifacts, and the C++ engine performs no
unit conversion at run time.

At each outer step all real degrees of freedom advance once under a single g-JF
force schedule (Eq.~\ref{eq:gjf}). Environment particles and membrane-contacting
protein carriers receive the production frictions derived in
Sec.~\ref{sec:timescale}. Zero-contact protein carriers retain the standard Upside
Ornstein--Uhlenbeck (OU) bath, whereas positive-g-JF-friction particles are
excluded from OU. Regenerated O and backbone-centre sites are not physical degrees
of freedom and receive neither update nor thermostat. Thus every real degree of
freedom has one stochastic bath at the same system temperature.

Preparation damping is deliberately separate from the production calibration.
Softened early stages use native dry-MARTINI damping, and full-core settling uses
strong FDT damping to dissipate energy released while packing overlaps are
resolved. Those stages prepare coordinates; their rate of motion has no physical
interpretation. The particle-mobility calibration is activated only in production.

The mechanical unit implied by $E_\mathrm{up}$, \AA, and $m_\mathrm{up}$ is
$0.20290$~ps. It is used only for unit conversion, not as the hybrid kinetic
clock. The declared protein clock is 40~ps per numerical step; the factor-four
MARTINI correction enters exactly once in the mobility calibration below. This
declared clock is a calibration convention, not evidence that a bonded lipid or a
membrane-coupled conformational transition evolves on physical time.
The two species couple through the interface interactions of
Sec.~\ref{sec:interface}, evaluated on the shared spline tables.

\subsection{The full-resolution dry-MARTINI membrane in Upside}
\label{sec:membrane}

Each DOPC lipid is its fourteen dry-MARTINI beads. Intramolecular connectivity is
taken from the dry-MARTINI topology: harmonic bonds (equilibrium lengths
$0.37$--$0.48$~nm, force constant $1250$~kJ\,mol$^{-1}$nm$^{-2}$) and harmonic
angles (force constants $25$--$45$~kJ\,mol$^{-1}$rad$^{-2}$), as Upside
\texttt{dist\_spring} and \texttt{angle\_spring} nodes. We adopt the
full-resolution lipid rather than a further-reduced single-particle lipid because
a reduced lipid cannot simultaneously reproduce packing, hydrophobic thickness,
and the orientational response that sets transmembrane tilt; an earlier
single-particle model (documented separately, now abandoned) failed there.

Nonbonded interactions---dry-MARTINI Lennard-Jones from the
\texttt{[nonbond\_params]} matrix, plus the Coulomb term for charged head-group
beads---are \emph{not} evaluated analytically at run time. Following
Upside's spline-table convention, every pair interaction is written to an HDF5
spline table (\texttt{combined\_energy\_grids}) ahead of the run and read back
through a single \texttt{martini\_potential} node with a cell-list neighbour
search. The hard $r^{-12}$ core is kept exactly as parameterized---no capping, no
softening---which is essential for correct bilayer packing and mechanics, for the
transport analysis of Sec.~\ref{sec:timescale}, and for confining an embedded
protein (Sec.~\ref{sec:interface}). The current Coulomb grid is truncated, not
shifted or tapered, at the 12~\AA{} neighbour cutoff. The resulting energy and
force discontinuity is an open numerical issue for charged BB--headgroup contacts
and must be resolved or shown irrelevant by timestep convergence before
quantitative protein--membrane equilibrium claims are made.

\subsection{Single-step g-JF Langevin lipid dynamics}
\label{sec:langevin}

We advance each lipid bead by one Gr\o{}nbech-Jensen--Farago (g-JF) Langevin step.
For bead $i$ of mass $m_i$, absolute friction $\alpha_i$, and
force $\mathbf{f}=-\nabla_i U$,
\begin{align}
b_i &= \Big(1+\tfrac{\alpha_i dt}{2m_i}\Big)^{-1}, \qquad
a_i = \Big(1-\tfrac{\alpha_i dt}{2m_i}\Big) b_i, \\
\mathbf{x}_i^{\,n+1} &= \mathbf{x}_i^{\,n} + b_i dt\,\mathbf{v}_i^{\,n}
  + \tfrac{b_i dt^2}{2m_i}\mathbf{f}_i^{\,n}
  + \tfrac{b_i dt}{2m_i}\boldsymbol{\beta}_i^{\,n+1}, \\
\mathbf{v}_i^{\,n+1} &= a_i\mathbf{v}_i^{\,n}
  + \tfrac{dt}{2m_i}\big(a_i\mathbf{f}_i^{\,n}+\mathbf{f}_i^{\,n+1}\big)
  + \tfrac{b_i}{m_i}\boldsymbol{\beta}_i^{\,n+1},
\label{eq:gjf}
\end{align}
with $\boldsymbol{\beta}_i^{\,n+1}=\sqrt{2\alpha_i k_BT\,dt}\,\boldsymbol{\xi}$ the
discrete thermal impulse ($\boldsymbol{\xi}$ standard Gaussian, the same
$\boldsymbol{\beta}$ in both updates). The shared force schedule evaluates
$\mathbf{f}^{\,n}$ before the update and $\mathbf{f}^{\,n+1}$ after moving the
g-JF particles. Bead
masses are the dry-MARTINI values. Positive-$\alpha$ particles are stochastic g-JF
particles and are excluded from the global OU bath; zero-$\alpha$ protein carriers
remain in that bath. g-JF has exact configurational statistics for a harmonic
oscillator within the Verlet stability limit, not unconditional exactness for an
arbitrary nonlinear force field. Static results therefore require timestep-halving
and distributional convergence checks. Its rate of motion is not generally
physical; that limitation is quantified in Sec.~\ref{sec:timescale}.

\subsection{Why molecular diffusion does not calibrate the one-step run}
\label{sec:designspace}

The preferred target is whole-DOPC lateral diffusion,
$D_\mathrm{DOPC}=11.5~\mu\mathrm{m}^2\,\mathrm{s}^{-1}$ after the MARTINI
factor-four correction. At 40~ps of declared protein time per numerical step,
this corresponds to a two-dimensional molecular-COM increment
\begin{equation}
  \left\langle \Delta r_{xy}^{2}\right\rangle
  =4D_\mathrm{DOPC}(40~\mathrm{ps})=0.184~\text{\AA}^2
  \quad\text{per step}.
\end{equation}
A one-step scan cannot reproduce it. Reducing friction through a free-draining
molecular estimate and nearly to zero raises the fitted molecular diffusion only
to approximately $0.175~\mu\mathrm{m}^2\,\mathrm{s}^{-1}$, while the
low-friction interface heats and loses hydrogen bonds. The molecular motion has
entered an inertial and collision-limited regime rather than the Einstein
overdamped regime assumed by a scalar-friction calibration.

The opposite error is also possible. A rejected mapping used
$\alpha=1666.67$ for a mass-6 environment bead, giving
$\alpha dt/(2m)=1.25$. Its momenta appeared thermal, but drift-removed DOPC COMs
moved only $0.17$--$0.24$~\AA{} over the nominal trajectory. Momentum temperature
therefore does not validate coordinate mobility. These two limits show why the
production friction is not presented as a molecular-diffusion fit.

\subsection{Bare-particle mobility and contact-local protein friction}
\label{sec:timescale}

Production instead calibrates a \emph{bare}, unbonded MARTINI-particle mobility.
Let $f_\mathrm{CG}=4$ be the MARTINI time correction,
$\Delta t_\mathrm{protein}=40$~ps per numerical step, and
$D_\mathrm{target}=11.5~\mu\mathrm{m}^2\,\mathrm{s}^{-1}$. The raw MARTINI
quantities are
\begin{equation}
  \Delta t_\mathrm{raw}=\frac{\Delta t_\mathrm{protein}}{f_\mathrm{CG}}
  =10~\mathrm{ps}, \qquad
  D_\mathrm{raw}=f_\mathrm{CG}D_\mathrm{target}
  =46~\mu\mathrm{m}^2\,\mathrm{s}^{-1}.
\end{equation}
The factor appears once: shortening the raw time and increasing the raw diffusion
are the two sides of the same mapping, and
$D_\mathrm{raw}\Delta t_\mathrm{raw}
=D_\mathrm{target}\Delta t_\mathrm{protein}$. Converting
$1~\mu\mathrm{m}^2\,\mathrm{s}^{-1}=10^{-4}$~\AA$^2$/ps and matching the
mean-square displacement accumulated over one numerical step gives
\begin{equation}
  D_\mathrm{bead,up}
  =\frac{D_\mathrm{raw}(10^{-4}~\text{\AA}^2/\mathrm{ps})
  \Delta t_\mathrm{raw}}{dt_\mathrm{up}}
  =5.11111~\text{\AA}^2/t_\mathrm{up},
  \label{eq:bead_diffusion}
\end{equation}
where $dt_\mathrm{up}=0.009$. The Einstein relation then sets the absolute
friction
\begin{equation}
  \alpha_\mathrm{bead}=\frac{k_BT}{D_\mathrm{bead,up}}.
  \label{eq:bead_friction}
\end{equation}
At the current $T=0.8\,T_\mathrm{up}$,
$\alpha_\mathrm{bead}=0.15652$ in Upside units. At any other temperature the
workflow recomputes Eq.~\ref{eq:bead_friction} from that same system temperature;
an independent bilayer calibration temperature is not allowed.

Every environment particle receives $\alpha_\mathrm{bead}$. For a protein
N/CA/C carrier $i$, the production model counts DOPC beads within the existing
12~\AA{} interaction range and assigns
\begin{equation}
  \alpha_i=n_i\alpha_\mathrm{bead}.
  \label{eq:contact_friction}
\end{equation}
Contact counts are recomputed from actual coordinates at stage handoff,
minimization promotion, and continuation, then held fixed within a trajectory
segment. Holding them fixed avoids silently introducing position-dependent
multiplicative noise. A zero-contact carrier remains on the ordinary Upside OU
bath; a positive-contact carrier uses g-JF and is excluded from OU.

Equation~\ref{eq:contact_friction} assumes independent additive bead drag. It
omits spatial and temporal friction correlations, anisotropy, hydrodynamic
coupling, and memory. It is therefore a reproducible local-damping model, not a
measurement of the friction experienced by a real membrane protein. Likewise,
Eq.~\ref{eq:bead_diffusion} is an input mobility for a hypothetical bare
particle. A bead bonded into DOPC shares the molecule's long-time COM diffusion,
so the trajectory must still report molecular diffusion independently. The H5
file stores the numerical step, system temperature, correction factor, raw and
corrected diffusion targets, bare-particle diffusion and friction, contact cutoff,
and per-carrier contact counts so this distinction is auditable.

\subsection{Membrane mechanics from a Monte-Carlo barostat}
\label{sec:mechanics}

Area compressibility $K_A$ and, through it, the bending modulus require sampling
the equilibrium area \emph{fluctuations}, $K_A=k_BT\langle A\rangle/\mathrm{Var}(A)$.
A weak-coupling (Berendsen- or damped-piston-style) barostat relaxes the box but
suppresses its fluctuations, so it cannot yield $K_A$. We therefore use a
\emph{Monte-Carlo} barostat: at fixed intervals, trial isotropic-in-plane and
normal box rescalings are applied by scaling lipid \emph{centre-of-mass} positions
(internal geometry preserved, so $\Delta U$ is intermolecular and reflects area
elasticity rather than bond stiffness) and accepted with the exact
constant-pressure Metropolis weight
$\exp\!\big[-\beta(\Delta U+P\Delta V)\big]\,(V'/V)^{N_\mathrm{mol}}$. Being a
Monte-Carlo move, it samples the correct isothermal--isobaric ensemble by
construction, independent of any coupling constant, and restores physical area
fluctuations. Consistent with the sub-diffusive lipid dynamics, the area
decorrelates slowly, so a converged $K_A$ requires long sampling; the barostat is
proposal and acceptance rule are verified, but quantitative convergence remains
open.

\subsection{The protein--membrane interface}
\label{sec:interface}

The protein keeps its native Upside force field. N, CA, and C are its physical
backbone carriers. O is regenerated from their geometry and the dry-MARTINI
backbone centre is constructed from N/CA/C/O. Neither derived site is integrated,
thermostatted, or counted in kinetic temperature. A dedicated coordinate node
constructs these sites during the forward pass. During the reverse pass the
BB-env sensitivity is copied with the dry-MARTINI backbone mass fractions
$14/54$, $12/54$, $12/54$, and $16/54$ to N, CA, C, and regenerated O,
respectively. Upside overwrites O during its integration cycle rather than
integrating it, so only the N/CA/C shares persist: $38/54$, or 70.37\%, of the
translational BB reaction. Differentiating the regenerated-O construction and
returning its share to N/CA/C would change this integration contract; it is not
done here. Raw N/CA/C/O--environment pairs are excluded so that the derived BB
interaction is counted exactly once.

\paragraph{SC-env (side chain $\leftrightarrow$ lipid).} Upside resolves a side
chain into one directional site, dry-MARTINI into several beads. This one genuine
resolution mismatch is bridged by a potential of mean force: for each residue
type, the multi-bead side chain is projected onto its single Upside site by a
Boltzmann average over the eliminated conformations,
\begin{equation}
U_\mathrm{SC\text{-}env}(\mathbf{r}) = -k_BT\,\ln\Big\langle
  \exp\!\big[-U_\mathrm{MARTINI}(\mathbf{r};\Omega)/k_BT\big]\Big\rangle_{\Omega},
\end{equation}
giving a smooth, orientation-resolved spline table (\texttt{sc\_table}). This is
the only PMF in the interface.

\paragraph{BB--environment.} The backbone is a single
site in both models, so nothing is integrated out: BB-env is the \emph{raw}
dry-MARTINI Lennard-Jones (plus Coulomb) interaction, from the same
\texttt{martini\_potential} spline as lipid--lipid. No softening, no PMF, no
capping. An early softened BB-env was tried and found to \emph{cause} the protein
to drift out of the bilayer (Sec.~\ref{sec:hybrid}); the raw field is retained.

\paragraph{Protein confinement.} No term keeps the protein in the membrane; its
equilibrium depth and orientation follow from the raw force field, as they should
from a correct force field rather than a restraint. There is no $z$-restraint and
no implicit-membrane solvation potential.

\subsection{Enhanced sampling and simulation protocol}
\label{sec:protocol}

Because both the intended observables and the slow intrinsic dynamics point the
same way, the model is intended for use with enhanced sampling. Temperature
replica exchange is supported natively (coordinate-swap convention: temperatures
are pinned to slots and the per-slot temperature is applied to the lipid
integrator, so each slot is a canonical ensemble); the lipid potential enters the
exchange criterion, and no momentum rescaling is needed. Well-tempered
metadynamics on a collective variable is available as a native bias node (a
history-dependent Gaussian bias read off a coordinate node, its gradient
propagated to the atoms), depositing on the simulation-round clock. These methods
can target equilibrium populations without preserving physical kinetics, but only
if the underlying finite-timestep distribution and the biased or replica-exchange
reweighting are validated. Availability of the algorithms does not establish
convergence of a particular trajectory.

Systems: a pure DOPC bilayer (structure, mechanics, transport), and transmembrane
proteins 1RKL (Ost4p) and 1AFO (glycophorin~A dimer). Lateral diffusion is
measured from unwrapped, mass-weighted 14-bead molecular COM trajectories after
bilayer COM drift removal; protein tilt is measured backbone-only.

The production BB-env and SC-env interactions are activated in a dedicated
\texttt{production\_handoff} stage. Throughout its minimization and burn-in the
complete protein is one rigid body: it may translate and rotate, but its internal
geometry cannot absorb resolution-handoff work. Only after burn-in is the stage
relabelled \texttt{production}, which removes the rigid-body constraint and
starts flexible sampling. No finite positional spring or gradual release is used,
and no preparation constraint contributes to the sampled production dynamics.

\section{Results}

Results are separated by the configuration that generated them. Baseline membrane
measurements predate the final particle-mobility mapping. Fresh 50,000-step hybrid
trajectories for 1RKL and 1AFO use $T=0.8$ for both the runtime thermostat and
friction calibration and store the expected
$D_\mathrm{bead,up}=5.11111$ and $\alpha_\mathrm{bead}=0.15652$. These runs test
the unified implementation; they do not by themselves establish timestep or
ensemble convergence.

\subsection{Baseline bilayer structure and mechanics}
\label{sec:structure}

The baseline bilayer remains assembled and gives area per lipid
$\approx65$~\AA$^2$ (exp.\ $\approx67$), hydrophobic
thickness (first-tail-bead separation) $\approx2.78$~nm (exp.\ $2.7$--$2.9$),
chain order parameter $\approx0.30$ (MARTINI range), and a fluid director
distribution (mean tilt $\approx22^\circ$, order $S\approx0.74$). The
phosphate--phosphate distance reads slightly thick ($\approx43$~\AA{}) but the
excess is head-group projection; the functionally decisive hydrophobic thickness,
to which a transmembrane protein matches, is on target. A zero-tension
Monte-Carlo-barostat run relaxes the area to a somewhat more condensed value than
the fixed-box state, indicating the force field is marginally over-cohesive---a
force-field fidelity note, reported and not tuned away. These values show that the
parameterization can form a bilayer, but they are not a converged validation of
the final temperature/friction/timestep combination. Quantitative $K_A$ and
bending-modulus claims remain open.

\subsection{Thermostat and transport checks}
\label{sec:barrier}

In the fresh 50,000-step trajectories, the last-200-frame protein kinetic-energy
means are 1.005 and 1.002 times $3k_BT/2$ for 1RKL and 1AFO; the corresponding
lipid ratios are 0.991 and 1.018.
This verifies that the implemented stochastic update thermalizes momenta at the
temperature it is given. It does not verify configurational sampling or a physical
clock. Multi-origin, drift-removed DOPC COM fits are strongly window-dependent:
$D=0.0053$--$0.0112~\mu\mathrm{m}^2\,\mathrm{s}^{-1}$ for 1RKL and
$0.0027$--$0.0103~\mu\mathrm{m}^2\,\mathrm{s}^{-1}$ for 1AFO over tested
20--1000~ns windows, versus the
$11.5~\mu\mathrm{m}^2\,\mathrm{s}^{-1}$ target. The window dependence itself
shows that a clean diffusive regime has not been reached, while every estimate is
more than three orders of magnitude below target. Molecular DOPC diffusion is
therefore demonstrably uncalibrated.

\subsection{Evidence boundary}
\label{sec:evidence}

\begin{table}[h]
\centering
\small
\begin{tabular}{
  >{\raggedright\arraybackslash}p{0.26\linewidth}
  >{\raggedright\arraybackslash}p{0.22\linewidth}
  >{\raggedright\arraybackslash}p{0.43\linewidth}}
\toprule
Quantity & Current status & Permitted interpretation \\
\midrule
One temperature and one bath per real degree of freedom &
Verified in workflow and H5 wiring &
The thermostat and friction calibration use the same $k_BT$ in the fresh H5 files. \\
Bare-particle friction of Eqs.~\ref{eq:bead_diffusion}--\ref{eq:bead_friction} &
Calibrated by construction &
An auditable effective mobility input; not the measured motion of a bonded bead. \\
DOPC molecular lateral diffusion &
Fails target &
Do not assign physical membrane time from the current trajectory. \\
Protein contact friction, Eq.~\ref{eq:contact_friction} &
Model assumption; unvalidated &
Useful as a local FDT-consistent damping rule, not as a physical protein
resistance or hydrodynamic tensor. \\
Bilayer and protein equilibrium populations &
Open &
Qualitative inspection only until timestep, cutoff, replica, and ensemble
convergence pass. \\
HDX analysis adapter &
Mechanically verified &
The data layout and EX2 calculation can be tested; present uptake and protection
factors are not quantitative predictions. \\
\bottomrule
\end{tabular}
\caption{The implemented calibration is narrower than a molecular kinetic
validation. An equilibrium observable is a possible target, not an automatic
guarantee that the available trajectory has sampled it.}
\label{tab:evidence}
\end{table}

\subsection{The raw force field confines the protein}
\label{sec:hybrid}

The raw field can keep a rigid 1RKL embedded, whereas replacing BB--environment
interactions with a smoothed PMF drives it out. The production interface therefore
retains the raw spline interactions and adds no confining potential. The current
virtual-coordinate node reconstructs O and the BB centre and applies the
historical partial mass-fraction route to persistent N/CA/C; direct table, type,
sign, cutoff, and pair-exclusion checks match the historical particle table.
These are implementation facts.

They do not establish secondary-structure stability. In the fresh, consistently
thermostatted $T=0.8$ run, DSSP assigns 13 helical residues to 1RKL at production
frame 0, 16 at the final frame, and as few as 5 during the trajectory; the
trajectory mean is 12.3. Thus substantial native helicity was already absent at
the production handoff and remains intermittent. Stage-resolved analysis found
23 helical residues throughout stages 6.0--6.6, but an unrestrained first
production-Hamiltonian minimization reduced that count to 9. The revised protocol
keeps the protein rigid through both interface minimization and burn-in. A reduced
end-to-end workflow test changed no persistent-backbone internal pair distance by
more than $4.6\times10^{-5}$~\AA{} before the flexible-production handoff. A new
full trajectory is still required, so the existing stage-7 result is not a test
of the revised protocol. By contrast, 1AFO has 54
helical residues initially, 52 finally, and a minimum of 51. Existing timestep
tests also show that $dt=0.009$ perturbs the coupled hard interface more strongly
than $dt=0.00225$.
In addition, charged pair splines are abruptly truncated at 12~\AA{}. A fresh
single-temperature run, a timestep series, and a cutoff-continuity test are
therefore required before interpreting helix populations thermodynamically. The
current trajectory can show that the workflow moves, remains finite, and keeps
the protein embedded over the observed window; it cannot establish a quantitative
helix-opening population. No interaction is softened or disabled.

\subsection{Application: hydrogen-deuterium exchange}
\label{sec:hdx}

In the common EX2 regime, conformational opening and closing equilibrate before
chemical exchange. The observed rate factorizes as
\begin{equation}
k_\mathrm{HDX} = \frac{k_\mathrm{op}}{k_\mathrm{cl}}\,k_\mathrm{chem}
= K_\mathrm{op}\,k_\mathrm{chem} = e^{-\Delta G_\mathrm{op}/k_BT}\,k_\mathrm{chem},
\end{equation}
so the individual opening and closing rates enter only through their equilibrium
ratio. The maintained analysis represents each frame by a binary structural
protection state, estimates the equilibrium protected probability
$p_\mathrm{protected}$ (with MBAR when replicas are available), and uses
\begin{equation}
  k_\mathrm{obs}=k_\mathrm{chem}
  \left(1-p_\mathrm{protected}\right)
  \label{eq:hdx_proxy}
\end{equation}
before calculating peptide uptake in experimental seconds. The MD time axis does
not appear in Eq.~\ref{eq:hdx_proxy}; the lipid-clock mismatch therefore does not
multiply or rescale the reported HDX seconds. It matters indirectly and
critically: slow or incorrect dynamics can leave
$p_\mathrm{protected}$ biased because the open and closed ensembles were not
sampled.

The hybrid adapter projects N/CA/C coordinates into the standard protein-only HDX
view while copying the full coupled-system potential, sampled temperature, time,
and compatible hydrogen-bond observables. It then reuses the existing protection,
MBAR, intrinsic-chemistry, and uptake workflow. Exact coordinate and observable
copying verifies this representation layer; it does not validate the source
ensemble or the protection model.

A quantitative membrane-protein HDX result requires all of the following:
\begin{enumerate}
  \item the requested experimental temperature is the single system temperature,
  and protein and lipid kinetic-temperature diagnostics agree with it;
  \item protection populations are stable to timestep reduction, trajectory
  blocking, and independent replicas or replica exchange;
  \item the coupled force field reproduces the relevant local helix-opening
  thermodynamics;
  \item an open amide is also water-accessible. Because dry-MARTINI has no water,
  accessibility must come from a calibrated local depth/water-activity model or
  an explicit-water reference.
\end{enumerate}
The fresh 1RKL source still fails these gates: it is a single trajectory with
strongly intermittent secondary structure, timestep convergence is unresolved,
the lipid ensemble has not reached a clean diffusive regime, and no calibrated
membrane-water accessibility is available. Its protection-state block convergence
and effective sample size must be recomputed before any HDX interpretation.
Current HDX outputs should therefore be used only to verify workflow wiring and
to perform clearly labelled qualitative sensitivity analyses. They should not be
reported as quantitative protection factors, residue opening free energies, or
peptide uptake predictions.

\subsection{Transmembrane helix tilt: a coupled-kinetics limit}
\label{sec:tilt}

In an earlier hybrid baseline, 1RKL stayed folded and embedded but its backbone
tilt remained near its $\sim40^\circ$ insertion value with only
$\pm0.2$--$0.4^\circ$ fluctuation, versus $\pm5$--$6^\circ$ in all-atom and
wet-MARTINI references that relax to $\sim10$--$22^\circ$. The frozen trajectory
proves inadequate relaxation but does not locate the hybrid model's equilibrium
tilt. Bidirectional starts or enhanced sampling of the tilt coordinate must
determine the free-energy surface before separating a kinetic limitation from a
force-field error.

\section{Discussion}

Three validation layers must remain separate. First, implementation validation
asks whether tables, forces, virtual-site derivatives, stochastic baths, and
metadata implement the stated equations. The one-temperature wiring, FDT impulse,
bare-particle friction, contact-count mapping, and HDX projection pass this layer.
Second, configurational validation asks whether the finite-timestep simulation
samples the intended equilibrium distribution and whether that distribution
matches the physical system. The present timestep dependence, charged-pair cutoff,
and short single-trajectory sampling leave this layer open. Third, kinetic
validation asks whether trajectory time corresponds to molecular time. The large
DOPC diffusion error and unvalidated protein drag fail this layer.

This hierarchy defines what can be trusted now. Input parameters and H5 metadata
can be trusted to describe the implemented stochastic model. Kinetic-temperature
averages can be trusted as momentum diagnostics. The bilayer remaining assembled,
the protein remaining embedded, and qualitative structural events can be used as
debugging observations tied to a particular run. One cannot yet trust absolute
DOPC or protein-coupled rates, a 40~ps physical interpretation of each trajectory
step, molecular diffusion, quantitative helix-opening populations, converged
bilayer elastic constants, or current HDX protection and uptake values.

Equilibrium methods remain the appropriate route forward, but equilibrium status
does not confer validity by itself. Timestep-halving, cutoff-continuity tests,
independent replicas, block convergence, and enhanced-sampling overlap must
establish the population before MBAR or EX2 postprocessing can make it
quantitative. Correcting a sampling defect by softening or disabling the physical
interface would instead change the model and is not acceptable.

\section{Outlook}

The principle---compose individually well-determined components through their
physical coupling, rather than train a bespoke potential for the combined
system---is not specific to lipids. Any system in which the parts are parameterized
but the \emph{combination} is data-poor is a candidate: a learned protein or
polymer model together with a physical model of an inorganic or engineered phase,
coupled by their raw cross-interactions. Nanoparticles are a concrete example---a
protein or membrane interacting with a nanoparticle whose surface chemistry is
physically describable but for which little protein-binding data exists. In every
such case, composing plausible components is only the starting point. Interface
forces, finite-step equilibrium, sampling convergence, and any kinetic mapping
must be validated independently.

\section{Conclusions}

A membrane protein and a full-resolution dry-MARTINI bilayer can be propagated
together through one g-JF force schedule with one temperature and one stochastic
bath per real degree of freedom. The factor-four mapping defines an explicit
bare-particle mobility, and the protein receives a reproducible contact-local
friction without modifying conservative SC--environment or BB--environment
interactions. This is an implementation result, not a physical-timescale
validation. Whole-DOPC diffusion is not matched, protein hydrodynamic friction is
not established, and current equilibrium and HDX populations have not passed
timestep and sampling gates. The method is therefore ready for controlled
validation and qualitative workflow studies, but current trajectories should not
be used for absolute membrane-coupled kinetics or quantitative HDX prediction.

\end{document}
